\appendix

\section{Appendix: operadic, PROPic and c/o structures}\label{appendixA}
% use *-form to suppress numbering



\subsection{The def\/inition of a c/o structure}
\label{coapp}
 Specify an object ${\mathcal O}(S,T)$ in some f\/ixed
symmetric monoidal category for each pair $S$ and $T$ of f\/inite
sets.  A {\it $G$-coloring} on ${\mathcal O}(S,T)$ is the further
specif\/ication of an object $G$ in this category and a morphism
$\mu : S\sqcup T \to {\rm Hom}({\mathcal O}(S,T),$G$)$, and we shall let
${\mathcal O}_\mu (S,T)$ denote this pair of data.

A $G$-colored ``closed/open'' or {\it c/o structure} is a
collection of such objects ${\mathcal O}(S,T)$ for each pair of
f\/inite sets $S$, $T$ together with a choice of weighting $\mu$ for
each object supporting the following four operations which are morphisms in the
category:


\vskip .1in

\noindent
\noindent{\it Closed gluing}:  $\forall\, s\in S$, $\forall\, s'\in S'$ with $\mu (s)=\mu '(s')$,
\[
\circ_{s,s'}: \  {\mathcal O}_\mu (S,T)\otimes{\mathcal O}_{\mu '} (S',T')\to {\mathcal O}_{ \mu ''}
(S\sqcup S'-\{ s,s'\}
, T\sqcup T');
\]


\noindent{\it Closed self-gluing}: $\forall \, s,s'\in S$ with $\mu (s)=\mu (s')$ and $s\neq s'$,
\[
\circ^{s,s'}: \ {\mathcal O}_\mu (S,T)\to {\mathcal O}_{\mu ''} (S-\{ s,s'\} , T);
\]



\noindent{\it Open gluing}:
$\forall  \, t\in T$, $\forall \, t'\in T'$ with $\mu (t)=\mu '(t')$,
\[
\bullet_{t,t'}: \ {\mathcal O}_\mu (S,T)\otimes{\mathcal O}_{\mu '} (S',T')\to {\mathcal O}_{\mu ''}
(S\sqcup S' , T\sqcup T'-\{t,t'\});
\]



\noindent{\it Open self-gluing}: $\forall \, t,t'\in T$ with $\mu (t)=\mu (t')$ and $t\neq t'$,
\[
\bullet^{t,t'}:  \ {\mathcal O}_\mu (S,T)\to {\mathcal O}_{\mu ''} (S , T-\{ t,t'\}).
\]


\noindent In each case, the coloring $\mu ''$ is  induced in the
target in the natural way by restriction, and we assume that $S\sqcup S'\sqcup T\sqcup T'-\{
s,s',t,t'\}\neq\varnothing$.

The  axioms are that the operations are equivariant
for bijections of sets and for bijections of pairs of sets, and the collection of all operations taken together
satisfy associativity.


Notice that we use the formalism of operads indexed by f\/inite
sets rather than by natural numbers as in \cite{MSS} for instance.


